\documentclass{article}

\usepackage[sglblindworkshop]{neurips_2026}

\usepackage[utf8]{inputenc}
\usepackage[T1]{fontenc}
\usepackage{hyperref}
\usepackage{url}
\usepackage{booktabs}
\usepackage{longtable}
\usepackage{amsmath,amsfonts,amssymb}
\usepackage{nicefrac}
\usepackage{microtype}
\usepackage{graphicx}

% Panel counts and every reported number, written by
% scripts/generate_results_tables.py, generate_arm_tables.py and
% layer_selection_audit.py from the stamped runs.
\input{generated/panel_description}
\input{generated/results_macros}
\input{generated/arm_macros}

\title{Only Two of Ten Nucleic Acid Foundation Models\\Encode Base-Pairing Partners}
\workshoptitle{XAI4Science: Knowledge Discovery and Trust through Interpretable Foundation Models}

\author{%
  Elliot Tower\\
  \texttt{elliot@elliottower.ai}\\
}

\begin{document}

\maketitle

\begin{abstract}
RNA foundation models are widely claimed to encode secondary structure, but
stems are GC-rich and loops are AU-rich, and no benchmark of these models
controls for it. We introduce a three-rung evaluation. Rung~1 asks whether a
model distinguishes stem from loop positions after controlling for base
composition; Rung~2 tightens the control to dinucleotide frequencies; Rung~3
asks whether a model knows which specific position pairs with which. We evaluate
\panelN{} Rfam~14.10 families spanning riboswitches, ribozymes, snRNAs, miRNA
precursors and CRISPR repeats, on ten models covering attention, state-space and
hybrid architectures, RNA and DNA pretraining, and 0.45M to 7B parameters. At
Rung~1, ERNIE-RNA discriminates stems from loops in \exceedErnieRNA{} of
\panelN{} families and RiNALMo in \exceedRiNALMo{}, where
\expectedFalsePositives{} are expected by chance; the best of the other eight
reaches nine. At Rung~3, only two of ten resolve which position pairs with
which: ERNIE-RNA at $0.888$ per-pair precision and RiNALMo at $0.872$, against
within-stem derangement chance rates of $0.124$ and $0.102$. The other eight sit
at chance. Their apparent excess comes entirely from choosing the read-out layer
per family. Three quarters of the two models' excess survives on synthetic
sequences carrying no evolutionary covariation. A randomly initialized ERNIE-RNA
clears every rung and reaches $0.490$, because re-initialization iterates
parameters and the architecture keeps a Watson--Crick table in a buffer; zeroing
it returns the control to its chance rate while the trained model keeps
$+0.725$.
\end{abstract}

\section{Introduction}
\label{sec:intro}

RNA foundation models are routinely described as encoding secondary structure,
on evidence from downstream task performance, linear probing, mutation
sensitivity, and attention--contact correlation. Each of those measurements
answers a different question and none controls for sequence composition. Deep
learning models for RNA structure prediction can fail to generalize across
families \citep{szikszai2022generalization}; whether learned representations
carry the same fragility has had less scrutiny.

\paragraph{The confound.} Stems and loops differ in composition. Across
\panelN{} Rfam families, stems average \panelStemGC{}\% G/C content against
\panelLoopGC{}\% in loops, and stem G/C exceeds loop G/C in \panelGCEnriched{}
of \panelGCScored{} families. Complement swaps at G/C positions produce larger
embedding perturbations than swaps at A/U positions for reasons unrelated to
pairing, so any stem-versus-loop contrast inherits the enrichment. A
representation can separate stems from loops while encoding nothing about which
position pairs with which.

\paragraph{Prior controls.} Composition-preserving null models are standard in
comparative RNA genomics, where dinucleotide-controlled shuffling was introduced
because mononucleotide shuffling understates false-positive rates for structure
prediction \citep{gesell2008dinucleotide, anandam2009multiperm}. No benchmark of
RNA language models applies one to model-derived structural signal. Current
evaluations score model output against reference structures, with a one-hot
sequence encoding as the closest control \citep{zablocki2025benchmarking}, which
asks whether an embedding beats raw nucleotide identity through a trained
predictor rather than whether the embedding's own structure sensitivity is
composition.

\paragraph{Graded resolution.} Structure awareness is not one property. We
separate three, each demanding more of a representation than the last:
(1)~stem-versus-loop discrimination, available to any embedding sensitive to
nucleotide composition; (2)~discrimination that survives nulls preserving
composition at one and two orders of sequence context; and (3)~partner
specificity, which requires that the representation encode pairing topology
rather than the nucleotides pairing implies.

\paragraph{Controls as the evidence.} A negative result at Rung~3 is worth
nothing if the test is weak, and a positive one is worth nothing if the baseline
is. We therefore put every baseline on trial before using it. The probe is run
on two constructed embeddings whose answers are fixed by design, one with partner
specificity and one without. The randomly initialized control is checked for
state that weight randomization does not reach. The chance rate is checked for
the selection the estimate receives. Two of the three needed repair, and the
repairs are priced.

\paragraph{Contributions.} We contribute (1)~a composition-controlled
perturbation evaluation of RNA base-pairing representation across ten models
spanning five architectures and three orders of magnitude in parameter count;
(2)~a probe validated against constructed embeddings with known answers;
(3)~evidence that a randomly initialized model clears every rung on an
architectural lookup table, with an ablation that prices it; and (4)~evidence
that per-instance layer selection accounts for the entire apparent partner
specificity of eight of the ten models.

\section{The Probe}
\label{sec:probe}

\paragraph{Perturbation specificity.} For a base pair $(i, j)$ we mutate
position $i$ to its Watson--Crick complement and record the change in every
position's embedding, $\delta_k = 1 - \cos(h_k, h_k')$, computed in float64.
Perturbation specificity is
$\mathrm{PS} = \delta_j - \max(\delta_{j-1}, \delta_{j+1})$: the partner must
move more than either of its own stem neighbors. Neighbors rather than arbitrary
positions make the comparison local, so a model that spreads perturbation along
a helix gains nothing. Per-pair precision is the fraction of eligible pairs for
which $\delta_j$ exceeds both neighbors.

\paragraph{The null.} Chance for this comparison is not a constant. Whether a
falsely assigned partner still outmoves its neighbors depends on how far a model
spreads a mutation: one whose perturbation decays sharply almost never ranks a
false partner first, and one that spreads it along the helix often does. We
permute partner assignments within each stem under a derangement, so every
position keeps its stem, its nucleotide and its layer while its assigned partner
is false by construction, and recompute per-pair precision. Measured this way the
rate runs from $0.000$ to $0.316$ across the ten models.

\paragraph{Panel.} We evaluate \panelN{} Rfam~14.10 families
\citep{kalvari2021rfam} spanning \panelClasses{}, with lengths
\panelLengthRange{}~nt (median \panelLengthMedian{}). Positions are labeled from
consensus dot-bracket annotation; pseudoknots, noncanonical pairs and gaps are
excluded. \rungThreeQualifying{} families carry at least 15 Watson--Crick pairs
with at least 3 per stem and qualify for Rung~3;
\rungThreeQuarantined{} pilot families are quarantined, leaving
\confirmatoryN{} confirmatory families and \chanceRateN{} carrying a derangement
chance rate. Full methods are in Appendix~\ref{app:methods}.

\section{Probe Validation}
\label{sec:validation}

Every arm of the panel is a model under test, so a model reading its own chance
rate admits two explanations the panel cannot separate: the representation lacks
partner specificity, or the probe cannot detect it. We therefore run the Rung~3
procedure unchanged on two constructed embeddings whose answers are fixed in
advance. In the oracle, each position's representation is a function of its own
nucleotide and its partner's. In the local control, it is a function of its own
nucleotide and its two sequence neighbors, and carries no partner information.

The oracle reaches per-pair precision $1.000$ against a derangement chance rate
of $0.099$. The local control reaches $0.155$ against $0.173$. Both use the same
statistic, the same null and the same code path as every model in
Table~\ref{tab:rung3}, over the \chanceRateN{} families of the Rung~3 analysis
set. The probe fires when partner specificity is present and reads chance when
it is absent. The two chance rates, $0.099$ and $0.173$ on systems whose answers
are known, are a second reason to measure chance per system rather than assume
it.

\section{Composition Absorbs Most Apparent Structure Signal}
\label{sec:composition}

Mutating a position to its Watson--Crick complement perturbs stem positions more
than loop positions in nine of ten models, with mean ratios from
$\ratioRNAFM{}$ (RNA-FM) down to near unity. Against a null that permutes stem
and loop labels within nucleotide strata, most of that signal is absorbed. At
$\alpha = 0.05$ per family, chance alone places \expectedFalsePositives{} of
\panelN{} families above the 95th percentile. ERNIE-RNA exceeds it in
\exceedErnieRNA{} families and RiNALMo in \exceedRiNALMo{}; RNA-FM, whose raw
ratio is the largest in the panel, exceeds it in \exceedRNAFM{}. Tightening the
null to preserve dinucleotide composition leaves ERNIE-RNA with
\dinucErnieRNA{} families and RiNALMo with \dinucRiNALMo{}, and every other
model with fewer than nine (Appendix~\ref{app:rungs12}). A raw ratio and a
composition-controlled count order the models differently, and neither
distinguishes a model that tracks pairing from one that tracks the nucleotides
pairing implies.

\section{Partner Specificity}
\label{sec:rung3}

\input{generated/rung3_table}

Mutating position $i$ perturbs its partner $j$ more than $j$'s own stem
neighbors in two models. RiNALMo reaches mean PS $=\psRiNALMo{}$ and ERNIE-RNA
$\psErnieRNA{}$, at per-pair precisions of $0.872$ and $0.888$ where each model's
own within-stem derangement puts chance at $0.102$ and $0.124$
(Table~\ref{tab:rung3}). The remaining eight exceed their own chance rates by at
most $0.122$, reached by Caduceus. Both positive models concentrate partner
specificity late: RiNALMo peaks at layer $\psPeakRiNALMo{}$ of
$\psLayersRiNALMo{}$ in $\psPeakCountRiNALMo{}$ families, ERNIE-RNA at layer
$\psPeakErnieRNA{}$ of $\psLayersErnieRNA{}$ in $\psPeakCountErnieRNA{}$.

Among the eight that do not clear their own chance rates, neither parameter
count nor pretraining domain orders anything: Evo reaches $\psEvo{}$ mean PS
with $75\times$ ERNIE-RNA's parameters, Caduceus at 7.7M outscores three of the
five RNA-pretrained models on mean PS, and all eight sit together at chance on
precision (Section~\ref{sec:selection}). The two that pass are both
RNA-pretrained, two of five against zero of five, on a comparison of ten models
that cannot resolve it (rank-biserial $\domainRankBiserial{}$,
$p = \domainP{}$).

\section{Layer Selection}
\label{sec:selection}

Per-pair precision is read at the layer where mean PS peaks, and that peak is
located separately in each of the \chanceRateN{} families. A model with a real
partner-specific response peaks in the same place every time, so choosing one
layer for the whole panel changes nothing. A model without one is taking the
maximum of a noise profile \chanceRateN{} times, against a chance rate that
receives no such choice.

\input{generated/layer_selection_table}

Two repairs give the same answer (Table~\ref{tab:layer-selection}). Choosing one
layer for the whole panel, and choosing the layer on half the families while
reading precision on the other half over $2{,}000$ splits, agree to three
decimals in every row. Under either, RiNALMo and ERNIE-RNA move by at most
$0.010$ and every other model in the panel falls to or below its own chance
rate. Caduceus goes from $0.340$ to $0.002$ and RNA-FM from $0.370$ to $0.126$,
against chance rates of $0.218$ and $0.261$. Depth does not predict how much the
per-family choice is worth---Evo has 32 layers and gains $0.041$, Caduceus has 17
and gains $0.340$, Spearman $+0.12$ over the thirteen arms ($p = 0.70$)---but the
presence of a peak does, which makes the gap between the two readings a test for
whether one exists.

The held-out estimate is selection-free by construction, and it is the number we
report for the two positive models: RiNALMo $0.882$ ($95\%$ interval over splits
$[0.822, 0.941]$) and ERNIE-RNA $0.885$ $([0.814, 0.955])$.

\section{The Randomly Initialized Control}
\label{sec:random-init}

Seven models were also run with every parameter re-initialized, holding the
architecture and the tokenizer and destroying the learned weights. ERNIE-RNA's
control clears each rung of the ladder. At Rung~1 it exceeds its
composition-preserving null in \exceedUntrainedErnieRNA{} of \panelN{} families
against a chance expectation of \expectedFalsePositives{}, retains
\retainUntrainedErnieRNA{} at Rung~2, and reaches per-pair precision $0.490$ at
Rung~3 against a chance rate of $0.241$.

\paragraph{A tensor randomization does not reach.} Re-initialization iterates a
model's parameters, and buffers are not parameters, so ERNIE-RNA's control keeps
the Watson--Crick table its architecture ships with:
\texttt{pairwise\_bias\_map}, a vocabulary-square buffer registered
non-persistent, whose only non-zero cells after the first forward pass are A--U
and U--A at $2.0$, C--G and G--C at $3.0$, and G--U and U--G at $0.8$. Any
architecture shipping a task-relevant lookup table outside
\texttt{named\_parameters()} produces a control that is not a null.

\input{generated/ablation_table}

\paragraph{Pricing it.} Ablating the table and the projection above it, with the
model's rebuild guard set so the first forward pass cannot restore it, separates
the two arms (Table~\ref{tab:ablation}). Writing excess as precision minus that
arm's own chance rate, the randomly initialized model falls from $+0.249$ to
$+0.012$ while the trained model falls from $+0.764$ to $+0.725$. The difference
between those losses is $+0.198$ ($95\%$ CI $[+0.069, +0.320]$) and the trained
model's remaining excess is $+0.725$ ($[+0.597, +0.836]$), both over
\chanceRateN{} families with \bootstrapB{} resamples. Predictions and criteria
were fixed before either ablated arm was computed.

\paragraph{The other controls.} Across the other nine models no buffer encodes
anything about base pairing; the rest carry position indices, rotary
frequencies, or all-zero token-type vectors, which survive randomization and
record where a nucleotide sits. Two controls nonetheless sat above their own
chance rates, RiNALMo by $0.105$ and RNA-FM by $0.093$, and both are layer
selection: read at a panel layer they fall to $-0.074$ and $-0.087$
(Table~\ref{tab:layer-selection}). Mean PS and per-pair precision disagree on
these arms. ERNIE-RNA's control has mean PS $3.4 \times 10^{-9}$ and per-pair
precision $0.490$, so a signed statistic at zero does not place a control at
chance.

\section{Covariation}
\label{sec:covariation}

Partner specificity could reflect sequence covariation rather than pairing
geometry. Compensatory mutations maintaining base pairing are overrepresented in
evolutionary alignments, and a model trained on them might predict that mutating
one position perturbs its covarying partner without encoding the physical
mechanism. Each panel family was therefore regenerated \synPerFamily{} times
with its dot-bracket structure preserved exactly and its nucleotides
reassigned---uniform Watson--Crick pairs at paired positions, uniform
elsewhere---giving \synN{} sequences that share structure with their natural
counterparts and carry no evolutionary covariation.

\input{generated/synthetic_table}

Mean PS falls by \synPSDropErnieRNA{} for ERNIE-RNA and \synPSDropRiNALMo{} for
RiNALMo. Excess precision over each arm's own chance rate falls by about a
quarter: ERNIE-RNA from \natExcessErnieRNA{} to \synExcessErnieRNA{} and RiNALMo
from \natExcessRiNALMo{} to \synExcessRiNALMo{}, retaining
\synRetainedErnieRNA{} and \synRetainedRiNALMo{}. The two statistics answer
different questions on these sequences. Mean PS is a magnitude, and synthetic
sequences have flatter composition, so every perturbation is smaller. Per-pair
precision asks whether a mutated position's partner moves more than that
partner's neighbors, which is what Rung~3 tests, and three quarters of it
survives on sequences with no covariation at all. Those sequences also destroy
natural sequence context, so the quarter that is lost is an upper bound on the
covariation contribution and the three quarters that remain are a lower bound on
the rest.

\section{A Composition-Biased Case Study}
\label{sec:htt}

The CAG repeat region in huntingtin exon~1 forms hairpins whose stability scales
with repeat count \citep{demezer2011mutant, krzyzosiak2012triplet}. We evaluated
five HTT exon~1 fragments (NM\_002111.7) with repeat counts 17, 21, 36, 40 and
60, folded with ViennaRNA \citep{lorenz2011viennarna}, through the panel's own
stages.

\input{generated/htt_table}

\httHighProbing{} of \httModels{} models reach linear-probe accuracy above
$0.97$ for stem-versus-loop classification, and \httHighProbingClean{} of those
exceed the composition-preserving null in no more than one of the five
fragments. The repeat region contains only C, A and G, so a probe has an extreme
composition bias available to it, and takes it.

Embedding distance from the wild-type fragment fails in the same way. It rises
monotonically with repeat count for all \httMonotone{} models with more than one
layer, at $\rho = +1.00$, and it already does so at the embedding layer, before
any block has run: the fragments differ in length and composition, and that is
enough. A monotone distance profile is not evidence that a representation
encodes the expansion. The one model whose final layer departs from the pattern
is RNA-FM, whose distances there collapse to $2 \times 10^{-7}$ in float64.

\section{Discussion}
\label{sec:discussion}

Composition-matched nulls absorb most of the apparent stem--loop sensitivity
across all ten models, and layer selection accounts for what remains in eight of
them. Two RNA-pretrained models retain partner-specific responses under a
selection-free read-out.

\paragraph{Two pathways.} RiNALMo and ERNIE-RNA reach partner specificity
differently. ERNIE-RNA injects a base-pairing attention bias, and that bias with
random weights above it resolves partners at $0.490$ against a chance rate of
$0.241$; ablating it removes the control's excess and costs the trained model
$0.039$. RiNALMo uses standard attention with no structural bias and reaches
comparable precision at $7.5\times$ the parameters. Whether the remainder
reflects a geometric internal model of folding or a statistical regularity of
paired positions is open, and answering it likely requires mechanistic analysis
rather than further behavioral benchmarking.

\paragraph{What a control has to be.} Both baseline failures here have the same
form: the control received a weaker version of the treatment the estimate
received. Weight randomization was assumed to remove everything the model was not
born with, and it missed a buffer. The estimate was a maximum over $L$ layers
and the null was not, so depth alone bought precision. Neither is visible in the
statistic being reported; both are visible in a baseline built to match. For
representation probes generally: enumerate \texttt{named\_buffers()} before
calling a randomly initialized model a null, and subject the null to whatever
selection the estimate receives.

\paragraph{For practice.} Any claim that an RNA representation encodes secondary
structure should report stem--loop sensitivity beside a nucleotide-stratified
null, so a reader can separate composition sensitivity from structure awareness.
Partner specificity is the more informative test: a model that perturbs stems
more than loops may be responding to G/C content, while one that selectively
perturbs the annotated pairing partner encodes positional pairing information.
Chance rates should be measured per model on a within-structure derangement
rather than assumed at $\nicefrac{1}{3}$: Caduceus at $0.340$ falls below
$\nicefrac{1}{3}$ while clearing its own rate by $0.122$, and the two
comparisons disagree about it.

\section{Limitations}
\label{sec:limitations}

The complement-swap metric tests sensitivity to single base-pair disruptions and
does not address tertiary contacts, pseudoknots, or long-range interactions. The
derangement null shares the model, the sequence and the layer with the estimate,
which is what makes it matched, and it therefore cannot detect a failure common
to all three. The synthetic control removes covariation and natural sequence
context together, so the quarter of the excess it costs bounds the covariation
contribution from above only. All structures are consensus secondary structures
from Rfam seed alignments rather than experimentally determined, and pretraining
overlap cannot be fully reconstructed. Ten models do not cover the model space,
and multi-nucleotide tokenization limits direct comparability for NT~v2 and
DNABERT-2.

\bibliographystyle{plainnat}
\bibliography{references}

\appendix
\section{Methods}
\label{app:methods}

\paragraph{Models.} Ten models spanning five architectures. RNA-pretrained:
RNA-FM (99M, transformer), RiNALMo (650M, transformer), ERNIE-RNA (86M,
transformer with a base-pairing attention bias), UTR-LM (1.2M,
transformer), SpliceBERT (19M, transformer). DNA-pretrained: NT~v2 (56M,
transformer, 6-mer), DNABERT-2 (117M, transformer, byte-pair), HyenaDNA (0.45M,
state-space), Caduceus (7.7M, state-space), Evo (7B, StripedHyena hybrid). Seven were
additionally run with every parameter re-initialized, xavier\_normal\_ on
matrices and normal\_$(0, 0.02)$ on vectors at a fixed seed.

\paragraph{Panel.} Rfam~14.10 seed-alignment consensus structures, minimum
20~nt with at least 5 stem and 5 loop positions. Pseudoknots, noncanonical
pairs and gap positions are excluded. \panelCurated{} families were curated and
\panelWithdrawn{} withdrawn on annotation review, leaving \panelN{}. Structures
are consensus secondary structures, not experimentally determined.

\paragraph{Perturbation.} For each annotated position we swap the nucleotide
with its Watson--Crick complement (A$\leftrightarrow$U, C$\leftrightarrow$G) and
measure cosine distance between wild-type and mutant embeddings, computed in
float64: at float32 the spacing of $1 - \cos$ near $1$ is $2^{-22}$, which
exceeds the signal for several models. For multi-nucleotide tokenizers,
token-level outputs are mapped to nucleotide positions through the tokenizer's
own offsets rather than a closed-form assumption about token width.

\paragraph{Nulls.} Rung~1 shuffles stem and loop labels within each nucleotide
type over 1{,}000 permutations, repeating layer selection inside every
permutation; Rung~2 conditions additionally on the 3$'$ neighbor. Rung~3
permutes partner assignments within each stem under a derangement, 500 per
family.

\paragraph{Aggregation.} The unit of replication is the RNA family. A family is
gate-qualifying for a model when its raw stem--loop ratio exceeds that model's
nucleotide-stratified null at the 95th percentile. Family-level bootstrap
resampling (\bootstrapB{} iterations, fixed seed) gives 95\% intervals.

\paragraph{Numerics and provenance.} Every cell pins TF32 off, cuDNN benchmark
off, and \texttt{torch.use\_deterministic\_algorithms(True)}, and records the
defaults it observed before pinning them. Each result file carries a stamp
naming the commit, the panel hash, the withdrawn families, the library
versions, the device and the weight condition; a run whose stamp differs from
the shards on disk discards them rather than resuming. Re-running Rung~3 under
this configuration reproduced every prior value byte for byte in twelve of
thirteen cells, the thirteenth differing at relative $10^{-6}$ in six families
because its earlier run predated the determinism pin.

\section{Rungs 1 and 2}
\label{app:rungs12}
\input{generated/rung1_table}

\input{generated/rung2_table}

\section{Supplementary tables}
\label{app:supplement}

\input{generated/panel_table}

\input{generated/attention_table}

\input{generated/transversion_table}

\input{generated/hypotheses_table}

\input{generated/provenance_table}

\end{document}
